\subsection{仿射空间 E 的乘积分解}

我们沿用命题 5 的记号。对于 \(0 \leqslant p \leqslant s\) ，设 \(\mathbf{E}_p\) 为群 \(\mathbf{T}_p'\) 在 \(\mathbf{E}\) 中诸轨道所成的集合，\(\varphi_p\) 为从 \(\mathbf{E}\) 到 \(\mathbf{E}_p\) 的典范映射。\(\mathbf{T}\) 在 \(\mathbf{E}\) 上的作用可过渡到商；特别地，\(\mathbf{T}_p\) 作用于 \(\mathbf{E}_p\) ，并且立即可见（例如在 \(\mathbf{E}\) 中取一个原点），\(\mathbf{E}_p\) 是一个仿射空间，以 \(\mathbf{T}_p\) 为其平移空间。置 \(\mathbf{E}' = \mathbf{E}_0 \times \cdots \times \mathbf{E}_s\) ；这是一个仿射空间，以 \(\mathbf{T} = \mathbf{T}_0 \oplus \cdots \oplus \mathbf{T}_s\) 为其平移空间。设 \(\varphi : \mathbf{E} \to \mathbf{E}'\) 为诸 \(\varphi_p\) 的乘积映射；由于 \(\varphi\) 与 \(\mathbf{T}\) 的作用交换，它是一个双射，甚至是仿射空间的同构。在下文中，我们把 \(\mathbf{E}\) 与 \(\mathbf{E}'\) 依照 \(\varphi\) 等同起来；映射 \(\varphi_p\) 于是等同于 \(\mathbf{E}'\) 到 \(\mathbf{E}_p\) 上的典范投影。

由于 W 保持 \(T_{p}^{\prime}\) 稳定，W 在 E 上的作用过渡到商，并定义 W 在 \(E_{p}\) 上的一个作用（对于 \(0 \leqslant p \leqslant s\) ），从而由限制得到 \(W_{p}\) 在 \(E_{p}\) 上的一个作用（对于 \(1 \leqslant p \leqslant s\) ）。另一方面，设 C 是一个室，\(i \in I\) ，并设 p 是使 \(i \in J_{p}\) 的整数。对于任意 \(x \in E\) ，我们有

\[
s _ {i} (\mathrm{C}) (x) = x - \lambda . e _ {i} (\mathrm{C}) \quad \text { with } \quad \lambda \in \mathbf {R}.
\]

由于 \(e_i \in \mathbf{T}_q'\) （对于 \(q \neq p\) ），由此得到

\[
\varphi_ {q} (w (x)) = \varphi_ {q} (x) \text {for} x \in \mathrm{E}, w \in \mathrm{W} _ {p}; 0 \leqslant q \leqslant s \text {and} q \neq p.
\]

因此，若 \(w = w_{1} \ldots w_{s} \in \mathrm{W}\) ，其中 \(w_{p} \in \mathrm{W}_{p}\) （对于 \(1 \leqslant p \leqslant s\) ），则

\[
w \left(\left(x _ {0}, \dots , x _ {s}\right)\right) = \left(x _ {0}, w _ {1} \left(x _ {1}\right), \dots , w _ {s} \left(x _ {s}\right)\right)\tag{5}
\]

对所有 \(x_{p} \in E_{p}\) （对于 \(0 \leqslant p \leqslant s\) ）成立。换言之，W 在 \(E = E'\) 上的作用恰是诸 \(W_{p}\) 在 \(E_{p}\) 上作用的乘积（我们置 \(W_{0} = \{1\}\) ）。由此可知，\(W_{p}\) 忠实地作用于 \(E_{p}\) ，并且 \(W_{p}\) 可等同于欧氏空间 \(E_{p}\) 的一个位移群（当然，平移空间 \(T_{p}\) —— 即 \(E_{p}\) 的平移空间 —— 赋予由 T 上的标量积所诱导的标量积）。

命题 6. (i) 群 \(W_{p}\) 是欧氏仿射空间 \(E_{p}\) 的一个位移群；它由反射生成；赋予离散拓扑后，它适当地作用于 \(E_{p}\) ；它是不可约的。

\begin{enumerate}
\def\labelenumi{(\roman{enumi})}
\setcounter{enumi}{1}
\tightlist
\item
  设 \(\mathfrak{H}_p\) 为超平面 \(\mathrm{H}\)（在 \(\mathrm{E}_p\) 中且满足 \(s_{\mathrm{H}} \in \mathrm{W}_p\) ）所成的集合。集合 \(\mathfrak{H}\) 由形如
\end{enumerate}

\[
\mathrm{H} = \mathrm{E} _ {0} \times \mathrm{E} _ {1} \times \dots \times \mathrm{E} _ {p - 1} \times \mathrm{H} _ {p} \times \mathrm{E} _ {p + 1} \times \dots \times \mathrm{E} _ {s},
\]

其中 \(p = 1,\dots ,s\) 且 \(\mathrm{H}_p\in \mathfrak{H}_p\)

\begin{enumerate}
\def\labelenumi{(\roman{enumi})}
\setcounter{enumi}{2}
\tightlist
\item
  每个室 C 都形如 \(E_{0} \times C_{1} \times \cdots \times C_{s}\) ，其中对每个 p，集合 \(C_{p}\) 都是 \(E_{p}\) 中由超平面集 \(H_{p}\) 所定义的一个室；此外，\(C_{p}\) 的墙是超平面 \(\varphi_{p}(H_{i}(C))\) ，其中 \(i \in J_{p}\) 。
\end{enumerate}

设 \(\mathbf{C}\) 是一个室。置 \(\mathrm{H}_i = \mathrm{H}_i(\mathrm{C})\) ，\(e_i = e_i(\mathrm{C})\) 且 \(s_i = s_i(\mathrm{C})\) （对于 \(i \in \mathbb{I}\) ，沿用第 4 号的记号）。

\begin{enumerate}
\def\labelenumi{(\roman{enumi})}
\tightlist
\item
  设 \(i\) 属于 \(\mathbf{J}_p\) ；由于 \(e_i \in \mathrm{T}_p\) 且 \(\mathrm{T}\) 是两两正交的子空间 \(\mathrm{T}_0, \mathrm{T}_1, \ldots, \mathrm{T}_s\) 的直和，\(\mathrm{T}\) 中与 \(e_i\) 正交的超平面形如 \(\mathrm{L}_i + \mathrm{T}_p'\) ，其中 \(\mathrm{L}_i\) 是 \(\mathrm{T}_p\) 中与 \(e_i\) 正交的超平面。仿射超平面 \(\mathrm{H}_i\) （\(\mathrm{E}\) 中的）形如 \(\mathrm{L}_i + \mathrm{T}_p' + x\) ，其中 \(x \in \mathrm{E}\) ，并且我们有
\end{enumerate}

\[
\mathrm{H} _ {i} = \mathrm{E} _ {0} \times \mathrm{E} _ {1} \times \dots \times \mathrm{E} _ {p - 1} \times \mathrm{H} _ {i} ^ {\prime} \times \mathrm{E} _ {p + 1} \times \dots \times \mathrm{E} _ {s}\tag{6}
\]

其中 \(H_{i}^{\prime}=L_{i}+\varphi_{p}(x)=\varphi_{p}(H_{i})\) 。现在立即可见，\(s_{i}\) 在 \(E_{p}\) 上的作用是超平面 \(H_{i}^{\prime}\) （\(E_{p}\) 中的）所对应的反射。于是，群 \(W_{p}\) 是 \(E_{p}\) 中由反射生成的位移群；适当性判别准则 (D'2) 的验证是直接的。最后，命题 5 (v) 表明 \(W_{p}\) 不可约。这就证明了 (i)。

\begin{enumerate}
\def\labelenumi{(\roman{enumi})}
\setcounter{enumi}{1}
\tightlist
\item
  由定理 1 的推论，集合 \(\mathfrak{H}_p\) 由形如 \(w_p(\mathrm{H}_i')\) 的超平面组成，其中 \(i\) 属于 \(\mathrm{J}_p\) 且 \(w_p\) 属于 \(\mathrm{W}_p\) 。进而，若 \(w = w_1 \ldots w_s\) 且 \(w_p \in \mathrm{W}_p\) （对所有 \(p\) ），则公式 (5) 与 (6) 蕴含
\end{enumerate}

\[
w (\mathrm{H} _ {i}) = \mathrm{E} _ {0} \times \mathrm{E} _ {1} \times \dots \times \mathrm{E} _ {p - 1} \times w _ {p} (\mathrm{H} _ {i} ^ {\prime}) \times \mathrm{E} _ {p + 1} \times \dots \times \mathrm{E} _ {s},\tag{7}
\]

由此立得 (ii)。

\begin{enumerate}
\def\labelenumi{(\roman{enumi})}
\setcounter{enumi}{2}
\tightlist
\item
  设 \(i\) 属于 \(\mathbf{J}_p\) ；由公式 (6)，开半空间 \(D_i\) （由 \(\mathbf{H}_i\) 界定且包含 \(C\) ）形如
\end{enumerate}

\[
\mathrm{D} _ {i} = \mathrm{E} _ {0} \times \mathrm{E} _ {1} \times \dots \times \mathrm{E} _ {p - 1} \times \mathrm{D} _ {i} ^ {\prime} \times \mathrm{E} _ {p + 1} \times \dots \times \mathrm{E} _ {s},
\]

其中 \(D_{i}^{\prime}\) 是一个由 \(H_{i}^{\prime}\) 界定的开半空间（在 \(E_{p}\) 中）。置 \(C_{p} = \bigcap_{i \in J_{p}} D_{i}^{\prime}\) ；由于 \(C = \bigcap_{i \in I} D_{i}\) ，立即可得

\[
\mathrm{C} = \mathrm{E} _ {0} \times \mathrm{C} _ {1} \times \dots \times \mathrm{C} _ {s};
\]

因此，诸集合 \(C_{p}\) 都非空，且由于 C 不与属于 H 的任何超平面相交，集合 \(C_{p}\) 不与属于 \(H_{p}\) 的任何超平面相交。于是 §1 第 3 号命题 5 表明，\(C_{p}\) 是由 \(H_{p}\) 在 \(E_{p}\) 中定义的一个室。利用 §1 第 2 号命题 4，容易看出 \(C_{p}\) 的墙是超平面 \(H_{i}^{\prime} = \varphi_{p}(H_{i})\) ，其中 \(i \in J_{p}\) 。

